\chapter*{Conclusion générale}
\addcontentsline{toc}{chapter}{Conclusion générale}
\label{chap:conclusion}

Tastify a permis de concevoir une application web dédiée à la gestion d'un restaurant. Le périmètre réalisé couvre les modules attendus pour ce type de projet : menu, tables, commandes, cuisine, stocks, ressources humaines, réservations, paiements, fidélité, avis clients, tableaux de bord et configuration.

Sur le plan technique, le projet repose sur Django REST Framework, MySQL, Redis, Celery, Django Channels, deux applications React et Docker Compose. Cette architecture répond à un besoin simple : faire circuler l'information entre les acteurs du restaurant sans mélanger les responsabilités. Le backend porte la logique métier, l'interface staff sert les rôles internes et le portail client garde son propre usage.

L'analyse de sentiment apporte une valeur utile au projet. Elle transforme les commentaires clients en indicateurs suivis par le gérant. Le rapport a présenté cette fonctionnalité en précisant les modèles utilisés, le secours local par mots-clés, les données exploitées, la conversion des labels et les limites de validation.

Le projet atteint un périmètre solide pour un PFA full-stack : les modules principaux sont reliés, les interfaces sont exploitables, les tests automatisés valident une large partie du socle et les choix techniques sont documentés. Les prochaines priorités relèvent donc de l'industrialisation : évaluation quantitative du sentiment, tests de charge exploitables, recommandation de plats plus avancée, prévisions de stock enrichies et déploiement de production renforcé.

Tastify constitue donc une solution cohérente et défendable dans le cadre d'un Projet de Fin d'Année. Il montre la capacité à concevoir un système full-stack réaliste, à relier plusieurs modules métier, à intégrer du temps réel et des tâches asynchrones, puis à documenter les choix techniques avec précision.
